\documentclass[11pt,twoside]{article}

\usepackage[
  paperwidth=8.5in, paperheight=11in,
  twoside,
  inner=1.05in, outer=2.35in,
  top=1.0in, bottom=1.05in,
  marginparwidth=1.75in, marginparsep=0.22in,
  headsep=16pt
]{geometry}

\usepackage{annotated}
\usepackage[colorlinks=true, allcolors=RuInk, linkcolor=RuInk]{hyperref}

\linespread{1.06}
\setlength{\parskip}{5pt}
\setlength{\parindent}{0pt}

\renewcommand{\leftmark}{Chapter 1 \;\textbullet\; The Real and Complex Number Systems}
\renewcommand{\rightmark}{The Annotated Rudin}

\begin{document}

% ============================================================
\thispagestyle{empty}
\vspace*{1.6in}
\begin{center}
{\sffamily\Huge\bfseries\color{RuInk} The Annotated Rudin}\\[10pt]
{\sffamily\large\color{RuGrey} A reader's companion to \emph{Principles of Mathematical Analysis}}\\[4pt]
{\sffamily\large\color{RuGrey} Third Edition}\\[40pt]
{\sffamily\normalsize\color{RuInk}\bfseries STYLE SAMPLE \textperiodcentered\ OPTION A \textperiodcentered\ B/W PRINT}\\[6pt]
{\sffamily\small\color{RuGrey} Chapter 1, \S\S1.1--1.21 --- roughly 8 of Rudin's pages}
\end{center}
\vfill
\begin{center}
\begin{minipage}{0.82\textwidth}
\small\color{RuGrey}
This sample exists so you can judge the \emph{style} before I commit to eleven
chapters. The mathematics is Rudin's; the boxes around it are the annotation
layer. Everything about the layout --- colours, box types, how much of Rudin's
own text gets quoted, how deep the commentary goes --- is still open.
\end{minipage}
\end{center}
\newpage

% ============================================================
\section*{How to read this book}
\addcontentsline{toc}{section}{How to read this book}

Rudin is terse by design. He states things once, in their sharpest form, and
leaves the reader to supply motivation, examples, and the missing algebra. This
companion supplies exactly those three things and nothing else. It is meant to
be read \emph{beside} Rudin, not instead of him.

This book is set for black-and-white printing, so the boxes are told apart by
their \emph{frames}, not by colour. Five kinds appear throughout.

\begin{roadmap}[What this chapter is for]
A fully boxed panel with a solid black title bar opens every chapter. It says
what the chapter is trying to achieve, gives the logical spine as a diagram, and
marks which parts you can safely skip. A matching panel closes the chapter with
the same ground shaded grey.
\end{roadmap}

\begin{rudinbox}{\textsc{2.27} Theorem}
\itshape A grey panel with a reverse-video number tab is Rudin, compressed.
Definitions and theorem statements appear here in his own phrasing (lightly
trimmed), so you always know what is his and what is mine. Proofs he gives are
reproduced when the annotation needs to point at specific lines.
\end{rudinbox}

\begin{unpack}
A \textbf{solid} rule down the left margin means unpacking. It answers ``what is
this actually saying,'' ``where did that formula come from,'' and ``why is the
definition phrased \emph{this} way and not the obvious way.'' If you are reading
Rudin for the first time, these are the boxes to read.
\end{unpack}

\begin{deeper}
A \textbf{dashed} rule and a faint grey ground means the box goes further than
the course requires: what breaks in a more general setting, which later theorem
this quietly sets up, what the standard generalisation looks like. Skippable on
a first pass.
\end{deeper}

\begin{pitfall}
Rules closing the box top and bottom, with a black label tab, mark a specific
wrong belief that students reliably form here --- and the counterexample that
kills it.
\end{pitfall}

\begin{proofwalk}[Reading the proof]
  \item Longer proofs get a numbered walk-through like this one, keyed to
        Rudin's lines. Each step says what is being done and, more importantly,
        \emph{why that step and not another}.
  \item Steps that look like rabbits pulled from hats --- Rudin's specialty ---
        get reverse-engineered, so you can see the choice being made rather
        than just verified.
\end{proofwalk}

\mnote{Margin notes carry one-line pointers: cross-references, the name of a
result, a warning to slow down.}
Finally, margin notes like the one beside this paragraph carry short pointers
that would interrupt the flow if inlined.

\newpage

% ============================================================
\section*{Chapter 1 \quad The Real and Complex Number Systems}

\begin{roadmap}[What this chapter is for]
Analysis is the study of limits, and a limit is a promise that some process
\emph{arrives} somewhere. The rationals break that promise: you can write down a
sequence of rationals closing in on $\sqrt 2$ with nothing at the end of it.
Chapter 1 fixes this once, by isolating the single property $\Q$ lacks --- the
least-upper-bound property --- and asserting a field that has it.

\medskip
The logical spine is short:

\medskip
\centering
\begin{tikzpicture}[scale=0.92, transform shape, node distance=5mm]
\node[rmbox, text width=23mm] (a) {\textbf{1.1}\\$\Q$ has gaps};
\node[rmbox, text width=23mm, right=of a] (b) {\textbf{1.8--1.10}\\$\sup$, $\inf$,\\the \lub{} property};
\node[rmbox, text width=23mm, right=of b] (c) {\textbf{1.19}\\$\R$ exists\\and has it};
\node[rmbox, text width=23mm, right=of c] (d) {\textbf{1.20, 1.21}\\Archimedes;\\$n$th roots};
\draw[rmarrow] (a) -- (b); \draw[rmarrow] (b) -- (c); \draw[rmarrow] (c) -- (d);
\end{tikzpicture}

\medskip
\raggedright
Everything after \rref{1.21} --- the extended reals, $\C$, Euclidean space --- is
setting up notation for later chapters. The appendix builds $\R$ from $\Q$ by
Dedekind cuts, discharging the promissory note in \rref{1.19}; you can read it
whenever you like, or never, without affecting anything else.
\end{roadmap}

% ------------------------------------------------------------
\subsection*{1.1 \; The gap in $\Q$}

Rudin opens by showing there is no rational $p$ with $p^2 = 2$ --- the familiar
parity argument --- and then does something less familiar. He sets
\[
  A = \{\, p \in \Q : p > 0,\; p^2 < 2 \,\},
  \qquad
  B = \{\, p \in \Q : p > 0,\; p^2 > 2 \,\},
\]
and shows \emph{$A$ has no largest member and $B$ has no smallest}. The tool is a
formula that arrives with no explanation:

\begin{rudinbox}{\textsc{1.1} Example (the key construction)}
\upshape
To each rational $p>0$ associate
\begin{align}
  q &\;=\; p - \frac{p^2-2}{p+2} \;=\; \frac{2p+2}{p+2}, \tag{1.3}\\[2pt]
  q^2 - 2 &\;=\; \frac{2\,(p^2-2)}{(p+2)^2}. \tag{1.4}
\end{align}
If $p \in A$ then $p^2 - 2 < 0$, so (1.3) gives $q > p$ and (1.4) gives $q^2 < 2$;
thus $q \in A$. If $p \in B$ then $p^2 - 2 > 0$, so $0 < q < p$ and $q^2 > 2$;
thus $q \in B$.
\end{rudinbox}

\begin{unpack}[Where does (1.3) come from?]
It is not magic and it is not Newton's method. Ask instead: \emph{which} map
$p \mapsto q$ would have $\sqrt2$ as a fixed point? Try the simplest shape, a
ratio of linear terms $q = (ap+b)/(cp+d)$. Demanding $q = p$ at $p = \sqrt 2$
means $p^2 = 2$ must fall out of $p(cp+d) = ap+b$, and $q = (2p+2)/(p+2)$ is the
smallest-integer solution: it gives $p^2 + 2p = 2p + 2$, i.e.\ $p^2 = 2$.

So (1.3) is the natural rational map pointing at $\sqrt 2$. The reason it moves
$p$ \emph{toward} $\sqrt2$ without ever crossing it is visible if you factor the
error. A short computation gives
\[
  q - \sqrt2 \;=\; \frac{\sqrt2\,(\sqrt2-1)}{p+2}\,\bigl(p - \sqrt2\bigr)
  \;=\; \frac{2-\sqrt2}{p+2}\,\bigl(p-\sqrt2\bigr).
\]
For $p > 0$ the multiplier $(2-\sqrt2)/(p+2)$ lies strictly between $0$ and
$(2-\sqrt2)/2 \approx 0.29$. Being \emph{positive} is what keeps $q$ on the same
side of $\sqrt2$ as $p$ --- that is Rudin's ``$q \in A$'' and ``$q \in B$.'' Being
\emph{less than one} is what makes the error shrink, which is why $q$ is strictly
better than $p$ and no member of $A$ can be largest.
\end{unpack}

\begin{deeper}[This is a contraction, eight chapters early]
The factorisation above says the map $p \mapsto (2p+2)/(p+2)$ contracts distance
to $\sqrt2$ by a factor of at most $0.29$ at every step. That is precisely the
hypothesis of the \emph{contraction principle}, \rref{9.23}, which Rudin uses to
prove the inverse function theorem. Chapter 1 quietly runs the argument by hand
before the general theorem exists.

Iterating from $p_0 = 1$ gives $1,\ \tfrac43,\ \tfrac{10}{7},\ \tfrac{24}{17},\
\tfrac{58}{41}, \dots$ --- the convergents of the continued fraction
$\sqrt2 = [1;2,2,2,\dots]$, alternating below and above only because we started
below. Compare Newton's method on $f(p) = p^2-2$, which gives $p \mapsto
(p^2+2)/(2p)$: faster (the error squares), but it always lands \emph{above}
$\sqrt2$, so it would prove ``$B$ has no smallest'' and say nothing about $A$.
Rudin's choice is the one that does both halves at once.
\end{deeper}

\begin{pitfall}[No largest member is not openness]
Openness is not the point and would not be enough. The content of \rref{1.1} is
that $A$ is bounded above in $\Q$ --- every element of $B$ is an upper bound ---
yet among those upper bounds there is no \emph{least} one, because $B$ has no
smallest element. A set can perfectly well have no largest member and still have
a supremum: $\{\,r \in \Q : r < 0\,\}$ has none, but its supremum $0$ exists in
$\Q$. The failure in \rref{1.1} is about the \emph{upper bounds}, not about $A$.
\end{pitfall}

% ------------------------------------------------------------
\subsection*{1.8--1.10 \; Supremum, infimum, and the property that matters}

\mnote{Read \rref{1.8} as two clauses: $\alpha$ \emph{is} an upper bound, and
nothing smaller is.}

\begin{rudinbox}{\textsc{1.8} Definition}
Let $S$ be an ordered set and $E \subset S$ bounded above. Suppose there is
$\alpha \in S$ with
\begin{enumerate}[label=(\roman*), leftmargin=2em, itemsep=1pt]
  \item $\alpha$ is an upper bound of $E$;
  \item if $\gamma < \alpha$ then $\gamma$ is not an upper bound of $E$.
\end{enumerate}
Then $\alpha$ is the \emph{least upper bound}, or \emph{supremum}, of $E$;
we write $\alpha = \sup E$. The infimum is defined symmetrically.
\end{rudinbox}

\begin{unpack}[Why not just say ``the largest element''?]
Because it usually does not exist. $E_1 = \{r \in \Q : r < 0\}$ has no largest
element, but the set of its upper bounds --- everything ${}\ge 0$ --- does have a
smallest, namely $0$. The supremum relocates the question from $E$, where the
answer is often ``none,'' to the set of upper bounds of $E$, where it is often
``yes.''

Clause (ii) is the whole definition. It is stated in contrapositive form, and
the version you will actually use in proofs is:
\[
  \text{for every } \varepsilon > 0 \text{ there is } x \in E
  \text{ with } x > \alpha - \varepsilon.
\]
Any $\gamma < \alpha$ fails to be an upper bound, which means \emph{something in
$E$ sticks out above $\gamma$}. Reach for that sentence every time a proof
begins ``let $\alpha = \sup E$.''
\end{unpack}

\begin{rudinbox}{\textsc{1.10} Definition}
An ordered set $S$ has the \emph{least-upper-bound property} if: whenever
$E \subset S$ is nonempty and bounded above, $\sup E$ exists in $S$.
\end{rudinbox}

\begin{pitfall}[The two hypotheses are not decoration]
``Nonempty'' and ``bounded above'' are both load-bearing, and dropping either
breaks the definition rather than merely weakening it. In $\R$, every real number
is an upper bound of $\emptyset$, so $\emptyset$ has no least upper bound; and
$\Z \subset \R$ is unbounded above, so it has none either. Any proof that
produces a supremum owes you a check of both --- in \rref{1.21} below, Rudin
spends two full sentences establishing exactly these before he writes
$y = \sup E$.
\end{pitfall}

% ------------------------------------------------------------
\subsection*{1.11 \; Greatest lower bounds come for free}

\begin{rudinbox}{\textsc{1.11} Theorem}
Suppose $S$ is an ordered set with the least-upper-bound property, $B \subset S$
is nonempty and bounded below, and $L$ is the set of all lower bounds of $B$.
Then $\alpha = \sup L$ exists in $S$, and $\alpha = \inf B$.
\end{rudinbox}

\begin{unpack}[The shape of the trick]
We are handed a hypothesis about suprema and asked for an infimum, so the only
possible move is to find a set whose \emph{supremum} is the infimum we want. The
set of lower bounds $L$ is the natural candidate: it sits entirely below $B$, so
$B$ pushes down on it from above, and its top edge is exactly where $B$ begins.
The picture is worth more than the algebra:
\[
  \underbrace{\cdots\cdots\cdots}_{\textstyle L}
  \;\bullet\;
  \underbrace{\cdots\cdots\cdots}_{\textstyle B}
  \qquad\text{with } \bullet = \sup L = \inf B .
\]
The two sets meet at a single point, and the theorem says that point is both.
\end{unpack}

\begin{rproof}
Since $B$ is bounded below, $L$ is not empty. Since $L$ consists of exactly those
$y \in S$ satisfying $y \le x$ for every $x \in B$, every $x \in B$ is an upper
bound of $L$. Thus $L$ is bounded above, and the hypothesis on $S$ gives
$\alpha = \sup L$ in $S$.

If $\gamma < \alpha$ then $\gamma$ is not an upper bound of $L$, hence
$\gamma \notin B$. It follows that $\alpha \le x$ for every $x \in B$; thus
$\alpha \in L$.

If $\alpha < \beta$ then $\beta \notin L$, since $\alpha$ is an upper bound of
$L$. So $\alpha \in L$ but $\beta \notin L$ whenever $\beta > \alpha$; that is,
$\alpha$ is a lower bound of $B$ and no larger element is. Hence
$\alpha = \inf B$.
\end{rproof}

\begin{proofwalk}
  \item \textbf{$L \ne \emptyset$, $L$ bounded above.} Both hypotheses of the
        \lub{} property must be checked before invoking it. Nonempty is the
        assumption that $B$ is bounded below, restated. Bounded above is the
        observation that every single $x \in B$ is an upper bound for $L$ ---
        by the very definition of ``lower bound of $B$.'' This is the one genuine
        idea in the proof.
  \item \textbf{$\alpha \in L$, i.e.\ $\alpha$ is a lower bound of $B$.} Rudin
        argues contrapositively and compresses two steps into one line. Unpacked:
        take any $\gamma < \alpha$. Since $\alpha$ is the \emph{least} upper
        bound of $L$, $\gamma$ is not an upper bound of $L$, so some
        $y \in L$ has $y > \gamma$. But $y$ is a lower bound of $B$, so
        $\gamma < y \le x$ for all $x \in B$, and in particular $\gamma \notin B$.
        No element of $B$ is therefore less than $\alpha$, which is to say
        $\alpha \le x$ for all $x \in B$.
  \item \textbf{Nothing bigger works.} If $\beta > \alpha$ then $\beta \notin L$,
        because $\alpha$ is an upper bound of $L$. So $\beta$ is not a lower bound
        of $B$. Combined with step 2, $\alpha$ is a lower bound and no larger
        element is --- which is the definition of $\inf B$.
\end{proofwalk}

\begin{deeper}[Why this is not a triviality]
The mirror-image statement is false in general order-theoretic settings, and the
reason the proof works here is that $L$ and $B$ are \emph{complementary} in a
strong sense: the upper bounds of $L$ include all of $B$. Take the greatest
lower bound property as a hypothesis instead and the same argument runs
backwards, so the two properties are equivalent for ordered sets. This is why
Rudin never states a separate ``greatest-lower-bound property'' axiom for $\R$:
\rref{1.19} gives him one, and \rref{1.11} hands him the other. Every later use
of $\inf$ --- lower Riemann--Stieltjes sums in \rref{6.2}, $\liminf$ in
\rref{3.16} --- rests on this one theorem.
\end{deeper}

% ------------------------------------------------------------
\subsection*{1.21 \; $n$th roots, and the rabbit out of the hat}

This is the theorem that repays the whole chapter: it is the first time the
\lub{} property produces a number that $\Q$ could not.

\begin{rudinbox}{\textsc{1.21} Theorem}
For every real $x > 0$ and every integer $n > 0$ there is one and only one
positive real $y$ with $y^n = x$.
\end{rudinbox}

Rudin sets $E = \{\, t > 0 : t^n < x \,\}$, checks $E$ is nonempty and bounded
above, puts $y = \sup E$, and rules out $y^n < x$ and $y^n > x$ in turn. The
first case runs like this:

\begin{rudinbox}{\textsc{1.21} Proof, the case $y^n < x$}
\upshape
The identity $b^n - a^n = (b-a)(b^{n-1} + b^{n-2}a + \cdots + a^{n-1})$ yields
\begin{equation}\tag{$\ast$}
  b^n - a^n < (b-a)\,n\,b^{n-1} \qquad (0 < a < b).
\end{equation}
Assume $y^n < x$. Choose $h$ so that $0 < h < 1$ and
\[
  h < \frac{x - y^n}{n\,(y+1)^{n-1}} .
\]
Put $a = y$, $b = y+h$. Then
\[
  (y+h)^n - y^n < h\,n\,(y+h)^{n-1} < h\,n\,(y+1)^{n-1} < x - y^n ,
\]
so $(y+h)^n < x$ and $y + h \in E$. Since $y+h > y$, this contradicts the fact
that $y$ is an upper bound of $E$.
\end{rudinbox}

\begin{unpack}[Reverse-engineering the choice of $h$]
Nobody guesses that $h$. It is found by writing down the goal and solving
backwards; Rudin simply deletes the derivation. Here it is.

\smallskip
\textbf{Goal.} Find $h>0$ with $(y+h)^n < x$. Equivalently,
$(y+h)^n - y^n < x - y^n$, where the right-hand side is a fixed positive number
because we assumed $y^n < x$.

\smallskip
\textbf{Bound the left side.} We cannot compute $(y+h)^n - y^n$, but $(\ast)$
with $a = y$, $b = y+h$ bounds it: $(y+h)^n - y^n < h\,n\,(y+h)^{n-1}$. So it is
\emph{enough} to make $h\,n\,(y+h)^{n-1} \le x - y^n$.

\smallskip
\textbf{Remove $h$ from the exponent.} That inequality still has $h$ buried
inside $(y+h)^{n-1}$, which we cannot solve for. Kill it by agreeing in advance
that $h < 1$; then $(y+h)^{n-1} < (y+1)^{n-1}$, a constant. This is the only
reason the condition ``$0 < h < 1$'' appears --- it is bookkeeping, not
mathematics.

\smallskip
\textbf{Solve.} It now suffices that $h\,n\,(y+1)^{n-1} < x - y^n$, i.e.
\[
  h < \frac{x-y^n}{n\,(y+1)^{n-1}},
\]
which is Rudin's line. Both conditions on $h$ are upper bounds, so take $h$
below the smaller of the two and such an $h$ exists.
\end{unpack}

\begin{deeper}[What the proof is really doing]
Strip the algebra and the argument is: \emph{$t \mapsto t^n$ is continuous, so if
$y^n$ were strictly below $x$ we could nudge $y$ upward and stay below $x$.} The
displayed choice of $h$ is a hand-rolled continuity estimate, written out
because continuity is not available until \rref{4.4} --- and \rref{4.4}'s
machinery is not available until $\R$ exists, which is what we are in the middle
of establishing. The same argument reappears in clean form as an application of
the intermediate value theorem (\rref{4.23}) once Chapter 4 is in place.

Note also which side each case attacks. Assuming $y^n < x$ contradicts ``$y$ is
an \emph{upper bound}''; assuming $y^n > x$ contradicts ``$y$ is the
\emph{least} upper bound.'' Both halves of \rref{1.8} get used exactly once,
which is a good sign that the choice of $E$ was the right one.
\end{deeper}

\begin{pitfall}[Why $E$ needs $t>0$]
Without the positivity restriction, $E = \{t : t^n < x\}$ is unbounded below for
odd $n$ and the set is a different animal for even $n$. More to the point, the
uniqueness claim is false without it: for $n=2$, $x=4$, both $2$ and $-2$ are
roots. Rudin's ``one and only one \emph{positive}'' is doing real work.
\end{pitfall}

\vspace{6pt}
\begin{recap}[Chapter 1 in one paragraph]
$\Q$ fails to be a home for analysis because bounded sets can have no least
upper bound (\rref{1.1}). Name that failure (\rref{1.8}--\rref{1.10}), assert a
field free of it (\rref{1.19}), and the consequences follow immediately: infima
exist too (\rref{1.11}), $\N$ is unbounded and $\Q$ is dense (\rref{1.20}), and
every positive real has every root (\rref{1.21}). Nothing in the rest of the
book uses more about $\R$ than this. When a later proof says ``by completeness,''
it means \rref{1.19} and nothing else.
\end{recap}

\end{document}
